\section{記号表}\label{sec:notation}

本付録は本文で使われるすべての記号とその意味をまとめる。

\begin{longtable}{@{}l p{0.78\linewidth}@{}}
  \toprule
  記号 & 意味 \\
  \midrule
  $\Nodes$            & 探索木のノード集合 \\
  $\Tree$             & 探索木 $(\Nodes, E)$ \\
  $n\in\Nodes$        & 個別のノード \\
  $\children{n}$      & $n$ の子集合 \\
  $\parentof{n}$      & $n$ の親 \\
  $\depth{n}$         & $n$ の深さ \\
  $\statefn{n}$       & ノード状態 $\in\{\StateOpen,\StateEval,\StateDone,\StateFailed\}$ \\
  $\nodescore{n}$     & $n$ のスカラー報酬 \\
  $\axes{n}$          & 軸ごとの評価ベクトル \\
  $\aggregate$        & 軸 $\to$ スカラー集約器 \\
  $\mathrm{div}(n)$   & 多様性ボーナス (\cref{eq:diversity}) \\
  $\fallbackscore{n}$ & 決定的フォールバックランキング (\cref{eq:fallback}) \\
  $\Tmax,\Cmax$       & ステップ・コスト予算 \\
  $\Rubric$           & ルブリックの木 \\
  $\rubricitem{i}$    & 葉 $i$ のカテゴリ \\
  $\rubricweight{i}$  & 葉 $i$ の重み \\
  $\judgefn{i}$       & 葉 $i$ の判官 \\
  $\paper$            & 論文草稿 \\
  $\paperscore{\paper}$ & 論文集約スコア (\cref{eq:paperscore}) \\
  $\refine$           & 改稿演算子 (\cref{eq:refine}) \\
  $\ckpt$             & チェックポイントディレクトリ \\
  $\lineage{n}$       & $n$ で終端する系譜の連鎖 \\
  \bottomrule
\end{longtable}
